\clearpage
\VerdierPhysicalPageAnchor{0148}%
\begingroup
\small
\setlength{\parskip}{0pt}%
\setlength{\emergencystretch}{1.5em}%
\setlength{\abovedisplayskip}{3pt plus 1pt minus 1pt}%
\setlength{\belowdisplayskip}{3pt plus 1pt minus 1pt}%

\begingroup\itshape
\VerdierSemanticAnchor{ch2-2-3-5-a-ii}
\noindent\hspace{1.2em}ii) $B\in\operatorname{Ob}(\mathcal B)$、
$X\in\operatorname{Ob}(\mathcal A)$ である任意の射 $B\longrightarrow X$ は
次の可換図式に組み込まれる：
\VerdierSemanticAnchor{ch2-2-3-5-a-ii-diagram}
\begin{center}
 \begin{tikzcd}[column sep=4.2em,row sep=2.6em]
 B \arrow[r] \arrow[dr] & X & {B'\in\operatorname{Ob}(\mathcal A\cap\mathcal B)} \\
 & B' \arrow[u] &
 \end{tikzcd}\quad .
\end{center}
\VerdierSemanticAnchor{ch2-2-3-5-a-definition}
上の二つの同値な性質をもつ $\mathcal D$ の充満三角部分圏 $\mathcal A$ を
\emph{右 $\mathcal B$-局所化部分圏}という。
\VerdierSemanticAnchor{ch2-2-3-5-b}
\par\noindent b) $\mathcal A$ が右 $\mathcal B$-局所化ならば、標準関手
\VerdierSemanticAnchor{ch2-2-3-5-b-functor}
\[
 \mathcal A/(\mathcal A\cap\mathcal B)\longrightarrow\mathcal D/\mathcal B
\]
は充満忠実で、対象上単射である。
\VerdierSemanticAnchor{ch2-2-3-5-c}
\par\noindent c) 次の条件は同値である：
\VerdierSemanticAnchor{ch2-2-3-5-c-i}
\par\noindent\hspace{1.2em}i) $\mathcal A$ は右 $\mathcal B$-局所化である。
\VerdierSemanticAnchor{ch2-2-3-5-c-ii}
\par\noindent\hspace{1.2em}ii) 標準関手
$Q\colon\mathcal D/(\mathcal A\cap\mathcal B)\longrightarrow
\mathcal D/\mathcal B$ に対し、部分圏
$\mathcal A/(\mathcal A\cap\mathcal B)\longrightarrow
\mathcal D/(\mathcal A\cap\mathcal B)$ は $Q$-右自由な対象からなる。
\endgroup

\VerdierSemanticAnchor{ch2-2-3-5-proof-a}
a) の同値性は TRIII と乗法的系 $S_{\mathcal D}(\mathcal B)$、
$S_{\mathcal A}(\mathcal A\cap\mathcal B)$ の定義 (2.1.7) から直ちに従う。
\VerdierSemanticAnchor{ch2-2-3-5-proof-b-setup}
b) を示す。$X$ と $Y$ を $\mathcal A$ の二対象（したがって
$\mathcal A/(\mathcal A\cap\mathcal B)$ の対象）とする。
\VerdierSemanticAnchor{ch2-2-3-5-proof-b-injectivity}
まず準同型
\[
 \operatorname{Hom}_{\mathcal A/(\mathcal A\cap\mathcal B)}(X,Y)
 \longrightarrow\operatorname{Hom}_{\mathcal D/\mathcal B}(X,Y)
\]
が単射であることを示す。
\VerdierSemanticAnchor{ch2-2-3-5-proof-b-quotient-functors}
$Q_1\colon\mathcal A\longrightarrow\mathcal A/(\mathcal A\cap\mathcal B)$ と
$Q_2\colon\mathcal D\longrightarrow\mathcal D/\mathcal B$ を標準商関手とする。
\VerdierSemanticAnchor{ch2-2-3-5-proof-b-fraction}
任意の元 $m\in\operatorname{Hom}_{\mathcal A/(\mathcal A\cap\mathcal B)}(X,Y)$
は $m=Q_1(s)^{-1}Q_1(f)$、
$s\in S_{\mathcal A}(\mathcal A\cap\mathcal B)$ と書ける (2.2.2)。その
$\operatorname{Hom}_{\mathcal D/\mathcal B}(X,Y)$ における像は
$Q_2(s)^{-1}Q_2(f)$ である。
\VerdierSemanticAnchor{ch2-2-3-5-proof-b-zero-criterion}
(2.2.4, a)) により、この像が零となるための必要十分条件は次の可換図式が
存在することである：
\VerdierSemanticAnchor{ch2-2-3-5-proof-b-zero-diagram}
\begin{center}
 \begin{tikzcd}[column sep=3.7em,row sep=2.5em]
 & {} \arrow[r,"u"] & {} \\
 X \arrow[ur,"f"] && Y \arrow[ul,"s"] \arrow[u,"t"'] &
 {t\in S_{\mathcal D}(\mathcal B)}
 \end{tikzcd}
\end{center}
\endgroup
